\newcommand{\placeholder}[1]{\textcolor{red}{[#1]}}

\section{Experiments}
\label{sec:experiments}

We design experiments to answer three questions. First, does step-independent training collapse as predicted by our dual mechanism analysis, and does episode-level training avoid this collapse? Second, how does the window size $h$ modulate the severity of failure, and is there a credit assignment cliff below some critical threshold? Third, can we directly measure advantage sign misattribution during training to validate the theoretical predictions of Section~\ref{sec:m2_sign_misattribution}?

\subsection{Experimental Setup}
\label{sec:experimental_setup}

\paragraph{Environment.}
We evaluate on ALFWorld~\cite{shridhar2021alfworld}, a text-based household environment where an agent interacts with a simulated kitchen, bedroom, and living room through natural language commands. % TODO: add shridhar2021alfworld to .bib
ALFWorld provides six task types (pick, clean, heat, cool, examine, pick two) that require sequential reasoning over multiple interaction steps. The agent issues text commands (e.g., ``go to countertop 1'', ``take knife 2 from drawer 3'') and receives textual observations describing the environment state. Episodes terminate after task completion or a maximum of 50 interaction steps. The reward is binary: $R = 1$ for successful task completion and $R = 0$ for failure. This binary reward structure makes ALFWorld a particularly challenging testbed for credit assignment: the only learning signal arrives at episode termination, and intermediate steps receive no direct feedback.

ALFWorld is well-suited for studying the step-independent failure because its tasks exhibit strong sequential dependencies. A typical task requires the agent to locate an object (steps 1--5), pick it up (step 6), navigate to a target receptacle (steps 7--10), and place the object (step 11). With step-independent windowing at $h = 2$, the window containing the placement action has no access to the exploration steps that determined which object to pick up.

\paragraph{Base model.}
We use Qwen3-1.7B (instruct variant) as the policy model. Before RL training, we perform supervised fine-tuning (SFT) on 2546 expert demonstrations collected from ALFWorld, training for 2 epochs with learning rate $5 \times 10^{-6}$. The resulting SFT checkpoint achieves approximately 6.25\% success rate on a held-out 32-episode validation set, providing a minimal but non-trivial starting point for RL. This low baseline success rate is important for our analysis: at $p = 0.0625$ with group size $K = 8$, Proposition 2 (Section~\ref{sec:formal_probability}) predicts that approximately 84.5\% of gradient updates carry incorrect sign under step-independent training.

\paragraph{RL training configuration.}
We train using the RAGEN framework~\cite{ragen2025} with the GRPO algorithm. The shared hyperparameters across all conditions are:
\begin{itemize}
    \item Group size $K = 8$ (8 rollouts per task instance for advantage estimation)
    \item Learning rate: $1 \times 10^{-6}$
    \item KL penalty coefficient: $\beta = 0.01$ (low-variance KL)
    \item Discount factor: $\gamma = 0.95$ (applied in step-independent return computation via Equation~\ref{eq:step_return}; not applicable to episode-level mode)
    \item Invalid action penalty coefficient: 0.1
    \item Training data: 16 task instances per training step
    \item Training steps: 15 (early-stopped based on convergence)
\end{itemize}

\paragraph{Evaluation protocol.}
We evaluate on a fixed 32-episode validation set every 5 training steps. The primary metric is success rate (pass@1): the fraction of episodes where the agent completes the task within 50 steps. We also report pass@$k$ ($k \in \{2, 4, 8\}$), computed by sampling $k$ rollouts per episode and counting success if any rollout succeeds. As a diagnostic metric, we track \textit{action validity rate}: the fraction of agent outputs that are syntactically valid environment commands (as opposed to malformed text, repeated outputs, or empty responses). Action validity serves as a proxy for basic policy coherence, distinct from task-solving capability. All reported numbers are single-seed unless otherwise noted; multi-seed results (3 seeds per condition) are marked as pending where applicable.

\subsection{Window-Size Ablation}
\label{sec:window_ablation}

To disentangle the effects of information loss (Mechanism 1) and sign misattribution (Mechanism 2) across different context lengths, we conduct a systematic ablation over window size $h$.

\paragraph{Conditions.}
We vary the step-independent history length $h \in \{2, 5, 10, 20\}$ and compare against episode-level training (equivalent to $h = T$, the full trajectory). Each condition uses the \texttt{step\_independent} training mode with \texttt{step\_independent\_history\_length} $= h$, except the episode-level condition which processes the complete trajectory as a single GRPO sample.

\paragraph{Design rationale.}
The five conditions span the spectrum from minimal context ($h = 2$, producing $N = 25$ windows per episode at $T = 50$) to full context ($h = T$, one sample per episode). Theory predicts two observable effects. Mechanism 1 (information loss) produces a performance ceiling that varies with $h$: smaller windows restrict the representable strategy space more severely, so step-1 performance should decrease monotonically with $h$. Mechanism 2 (sign misattribution) produces a training-time degradation that is present for all $h < T$: as long as windows exist, the advantage sign is determined by episode outcome rather than local action quality, so performance should decline over training steps regardless of window size. The combined prediction is a credit assignment cliff: performance drops sharply below a critical $h$ value, with smaller $h$ exhibiting both lower initial performance and faster collapse.

\paragraph{Controls.}
All conditions share the same GRPO hyperparameters (group size, learning rate, KL penalty, discount factor). Each condition runs for 15 training steps with validation every 5 steps. We plan 3 seeds per $h$ value to assess variance. All $h$-sweep conditions are currently single-seed. Multi-seed replication (3 seeds per $h$) is planned.% Sources: exp028c ($h=2$), exp029 (episode-level).

\paragraph{Post-hoc evaluation.}
To isolate the training-layer effect of sign misattribution from the inference-layer effect of context truncation, we conduct full-context evaluation of checkpoints trained with truncated windows. Concretely, a model trained with $h = 2$ is evaluated using the complete interaction history at inference time. If the trained model performs better under full-context evaluation than under $h = 2$ evaluation, the difference reflects pure inference-time information loss (Mechanism 1). If it still performs worse than the episode-level trained model even under full-context evaluation, the gap reflects training-layer damage from sign misattribution (Mechanism 2).

\subsection{Results}
\label{sec:results}

\subsubsection{Episode-Level vs. Step-Independent Training}
\label{sec:results_main}

The central empirical finding is the contrast between episode-level and step-independent training. Episode-level GRPO achieves 28.125\% validation success at training step 15, a 4.5$\times$ improvement over the SFT baseline of approximately 6.25\%. Step-independent training with $h = 2$ collapses to 0\% validation success by step 10.

Table~\ref{tab:main_results} summarizes the central comparison. All seed results are reported individually due to high inter-seed variance under step-independent conditions.

\begin{table}[t]
\centering
\caption{1.7B ALFWorld results across credit assignment methods. OOD: held-out test tasks; In-Dist: 32-episode validation set. All seeds shown; no aggregation due to high inter-seed variance.}
\label{tab:main_results}
\small
\begin{tabular}{@{}llcc@{}}
\toprule
\textbf{Method} & \textbf{Seed} & \textbf{OOD (\%)} & \textbf{In-Dist (\%)} \\
\midrule
Step-indep.\ ($h{=}2$, windowed) & 1 & 0 & 0 \\  % Cell A
                                  & 2 & 3.73 & 9.375 \\  % Cell A seed2
\midrule
Step-indep.\ (factorial)          & 1 & 0.75 & 20 \\  % Cell B
                                  & 2 & 0.75 & 55.5 \\  % Cell B seed2
\midrule
Episode-level (full context)      & 1 & 23.88 & 28.125 \\  % Cell D, exp029
\bottomrule
\end{tabular}
\end{table}

\placeholder{fig6\_training\_curves: Training curves showing episode-level vs step-independent $h=2$. Dual y-axis: success rate (left) and action validity (right). X-axis: training steps 0--15.}

The collapse under step-independent training follows the two-phase pattern predicted by our analysis. In the first phase (steps 0--2), task success drops from 18.75\% to single digits while action validity remains high at approximately 89.3\%. The model retains basic command formatting but loses the ability to compose multi-step strategies. In the second phase (steps 2--4), action validity itself collapses from 89.3\% to 43.1\%. Accumulated wrong-sign gradient updates, which initially eroded only strategic capability, propagate to corrupt the action space representation. By step 10, the agent produces predominantly malformed outputs and achieves no task completions.

Episode-level training exhibits no such collapse. Action validity remains above 95\% throughout training, and success rate increases monotonically. The single advantage shared across all steps in an episode preserves correct actions even within failed episodes (since all steps receive the same sign, no individual step is singled out for incorrect suppression). Cross-step dependencies are maintained through the full interaction history, so the model can learn conditional strategies that require information from early exploration steps.

\subsubsection{Window-Size Ablation}
\label{sec:results_hsweep}

\begin{table}[t]
\centering
\caption{Window-size ablation (1.7B, single seed). OOD success rate at the final training step. $h = \text{full}$ denotes episode-level training (exp029).}
\label{tab:h_sweep}
\small
\begin{tabular}{@{}lc@{}}
\toprule
\textbf{Window size $h$} & \textbf{OOD Success (\%)} \\
\midrule
$h = 2$  & 0 \\
$h = 5$  & 0 \\
$h = 10$ & 6.25 \\
$h = 20$ & 15.6 \\
$h = \text{full}$ (episode-level) & 23.88 \\
\bottomrule
\end{tabular}
\end{table}

\placeholder{fig5\_h\_sweep\_curve: $h$ vs. validation success at step 15. Error bars from 3 seeds. Shows credit assignment cliff. Horizontal dashed line at SFT baseline (6.25\%).}

The ablation reveals a credit assignment cliff consistent with our theoretical predictions (Table~\ref{tab:h_sweep}). OOD success rate is 0\% for $h \leq 5$, rises to 6.25\% at $h = 10$ and 15.6\% at $h = 20$, and reaches 23.88\% under episode-level training ($h = T$). The transition between $h = 5$ and $h = 10$ marks the onset of non-trivial generalization. Results are single-seed; multi-seed replication is in progress.

The dual mechanism theory predicts three observable signatures in the $h$-sweep. First, step-1 performance should vary across $h$ values, reflecting different information-loss ceilings (Mechanism 1). Second, for all $h < T$, performance should decline over training steps as sign misattribution accumulates (Mechanism 2). Third, the rate of decline should be approximately constant across $h$ values (since the sign correlation from Equation~\ref{eq:sign_correlation} depends on the success rate $p$ and group size $K$, not on $h$ directly), but the onset of visible collapse should be earlier for smaller $h$ because the lower ceiling leaves less room before reaching zero.

From the single-seed data available, the two-phase collapse at $h = 2$ is confirmed. Action validity at step 2 is 89.3\%, dropping to 43.1\% by step 4. This pattern, strategic collapse preceding formatting collapse, is the empirical signature of sign misattribution operating on a model that initially retains valid action representations. The progressive nature distinguishes Mechanism 2 from Mechanism 1: information loss would produce a constant (low) performance from step 1, not a declining trajectory.

\subsubsection{Dense Reward Baseline}
\label{sec:results_dense}

A natural hypothesis is that step-independent training fails due to sparse reward, and that providing dense per-step rewards might resolve the credit assignment problem. We test this with a dense reward configuration that adds $+0.15$ bonus for each syntactically valid action, applied under step-independent training with $h = 2$.

The dense reward experiment (exp013f) produces a striking dissociation. Action validity increases from 60\% to 97.5\% over training. The per-step bonus successfully shapes low-level formatting behavior. However, task success remains at 0\% throughout training. The agent learns to produce perfectly formatted commands that satisfy the validity reward without solving any tasks.

This result confirms that the failure of step-independent training is not a sparse reward problem. Dense reward shaping creates an optimization target that the agent can maximize (action formatting) without addressing the underlying credit assignment pathology. The sign of the task-level advantage still corrupts multi-step strategy learning, while the dense reward merely teaches the agent to produce syntactically correct (but strategically meaningless) outputs. Reward shaping and sign misattribution operate on different dimensions of behavior and cannot compensate for each other.

\subsection{Empirical Verification of Sign Misattribution}
\label{sec:sign_verification}

The theoretical analysis of Section~\ref{sec:m2_sign_misattribution} predicts specific quantitative properties of the advantage signal under step-independent training. We design a direct measurement protocol to validate these predictions.

\paragraph{Methodology.}
During step-independent training, we log per-window advantage information for each training episode. For window $k$ in episode $i$, we record: (1) the advantage sign $\text{sign}(\hat{A}_k^{(i)})$, (2) the episode outcome $R_{\text{episode}}^{(i)} \in \{0, 1\}$, and (3) a local action quality proxy based on the fraction of valid actions within the window. We then compute three statistics across all training windows:
\begin{itemize}
    \item \textbf{Sign-outcome correlation}: the Pearson correlation between $\text{sign}(\hat{A}_k)$ and $R_{\text{episode}}$. Theory predicts this correlation approaches 1.0, since the advantage sign is determined almost entirely by whether the episode succeeded.
    \item \textbf{Sign-quality correlation}: the Pearson correlation between $\text{sign}(\hat{A}_k)$ and local action quality. Theory predicts this correlation approaches 0, since the advantage sign carries no information about whether the window's actions were locally correct.
    \item \textbf{Correct sign fraction}: $P(\text{correct sign} \mid \text{update})$, defined as the probability that a window with locally good actions receives a positive advantage and a window with locally bad actions receives a negative advantage. Proposition 2 (Section~\ref{sec:formal_probability}) predicts this equals approximately 15.5\% at $p = 6.25\%$, $K = 8$.
\end{itemize}

\paragraph{Predictions.}
The sign misattribution mechanism makes three testable predictions:
\begin{enumerate}
    \item Sign-outcome correlation $\approx 1.0$ (advantage sign is a proxy for episode success, not local quality)
    \item Sign-quality correlation $\approx 0$ (advantage sign is independent of whether local actions were good)
    \item $P(\text{correct sign}) \approx 15.5\%$ (derived from Proposition 2 at the observed success rate)
\end{enumerate}
If all three predictions are confirmed empirically, the sign misattribution mechanism is validated as the causal driver of training collapse, rather than alternative explanations such as reward sparsity, KL divergence, or optimization instability.

Table~\ref{tab:sign_verification} reports the empirical measurements. The credit agreement rate---the fraction of windows where the advantage sign matches local action quality---is 52.07\%, indistinguishable from the 50\% expected under random sign assignment. The point biserial correlation between advantage sign and local action quality is $r = 0.068$, confirming near-zero dependence. Additionally, a discount factor sweep reveals that \textit{all} $\gamma$ values produce collapse, including $\gamma = 1.0$ (no discounting). The $\gamma = 1.0$ result rules out temporal discounting as the failure cause and isolates sign misattribution as the root mechanism.

\begin{table}[t]
\centering
\caption{Empirical verification of sign misattribution under step-independent training ($h{=}2$). Credit agreement rate measures the fraction of windows where advantage sign matches local action quality; 50\% indicates random assignment.}
\label{tab:sign_verification}
\small
\begin{tabular}{@{}lcc@{}}
\toprule
\textbf{Metric} & \textbf{Measured} & \textbf{Predicted} \\
\midrule
Credit agreement rate & 52.07\% & $\approx$ 50\% (random) \\
Point biserial $r$ (sign vs.\ quality) & 0.068 & $\approx$ 0 \\
Collapse at $\gamma = 1.0$ (no discounting) & Yes & Yes \\
\bottomrule
\end{tabular}
\end{table}

\subsection{Additional Baselines}
\label{sec:additional_baselines}

\paragraph{GAE baselines.}
To distinguish the effect of episode-level optimization from GRPO-specific design choices (group normalization, no value function), we compare against Generalized Advantage Estimation (GAE)~\cite{schulman2016gae} under two configurations. GAE($\lambda = 0.95$, $\gamma = 0.95$) provides standard TD-based advantage estimation with a learned value function, representing the conventional actor-critic approach to temporal credit assignment. GAE($\lambda = 1$, $\gamma = 1$) computes Monte Carlo returns equivalent to episode-level GRPO without group normalization, isolating the contribution of group-relative advantage computation. \placeholder{plan\_003b: GAE baseline execution. Columns: Method, Val Success (\%), Action Valid (\%), Training Stability (variance across steps).}

\paragraph{Selective verification.}
Building on the episode-level trained agent (Section~\ref{sec:selective_verification}), we train a selective verification policy that decides when self-checking is informative. This component is orthogonal to the credit assignment analysis and demonstrates a downstream application of the competent base agent produced by episode-level training. \placeholder{tab4\_selective\_comparison: Selective verification results. Columns: Method (Always-Verify, Never-Verify, Selective), Success (\%), Verify Rate (\%), Cost Reduction (\%). Pending Phase B experiments.}

\subsection{Cross-Scale Verification}
\label{sec:cross_scale}

The 1.7B experiments establish that step-independent training collapses while episode-level training succeeds. Both mechanisms could in principle be scale-dependent: a larger model might compensate for information loss through greater capacity, or resist sign misattribution through more stable optimization. We test scale invariance by replicating the key factorial conditions on Qwen3-8B.

\paragraph{Setup.}
The 8B SFT baseline uses the same training protocol (2{,}546 demonstrations, 2 epochs, learning rate $5 \times 10^{-6}$). It achieves 68.75\% in-distribution success ($T{=}0$, deterministic) and 80.6\% OOD success ($T{=}0.4$, stochastic sampling), far above the 1.7B baseline of 12.5\%.% exp_id: 8b_det_eval_full (in-dist), 7b_sft_eval_ood (OOD)
This higher starting point provides a stronger test: improvement over an already-competent baseline would confirm that episode-level RL adds value beyond what scale alone provides.

\paragraph{Information loss.}
Restricting the 8B model's observation window to $h{=}2$ at inference time drops in-distribution success from 68.75\% to 31.25\%, a 37.5 percentage point gap.% exp_id: 8b_info_loss_eval
The relative magnitude is comparable to the 1.7B case. Larger models do not implicitly reconstruct missing context from a 2-step window.

\paragraph{Cell A collapse.}
Step-independent training on the 8B model collapses within two gradient steps: in-distribution success drops from 28.1\% (before training) to 12.5\% (step~1) to 1.5\% (step~2).% exp_id: 8b_cell_a_episode
The collapse trajectory matches the 1.7B pattern, consistent with the sign misattribution rate predicted by Proposition~\ref{prop:correct_sign}.

\paragraph{Cell D: episode-level training.}
Episode-level GRPO reaches 96.9\% in-distribution success at step~10 and 90.3\% OOD success (134 tasks, $T{=}0.4$), a 9.7 percentage point improvement over the SFT baseline on OOD tasks.% exp_id: 8b_cell_c_episode_v9d, 8b_v9d_ood_eval_v5
Training reward saturates to 100\% at step~12. The gain over an 80.6\% SFT baseline confirms that correct credit assignment enables RL to improve beyond supervised learning even from a strong initialization.

Table~\ref{tab:8b_results} summarizes the 8B conditions.

\begin{table}[t]
\centering
\caption{8B model results on ALFWorld. In-Dist: 32-task validation set; OOD: 134-task held-out test set. SFT in-distribution and SFT $h{=}2$ use deterministic decoding ($T{=}0$); all OOD results and RL evaluations use $T{=}0.4$. Cell~A reports the value at step~2 (post-collapse). Cell~D reports the best checkpoint (step~10).}
\label{tab:8b_results}
\small
\begin{tabular}{@{}lcc@{}}
\toprule
\textbf{Method} & \textbf{In-Dist (\%)} & \textbf{OOD (\%)} \\
\midrule
SFT baseline & 68.75 & 80.60 \\  % exp_id: 8b_det_eval_full (in-dist T=0), 7b_sft_eval_ood (OOD T=0.4)
SFT $h{=}2$ & 31.25 & --- \\  % exp_id: 8b_det_eval_h2 (T=0)
Cell A: step-indep.\ ($h{=}2$) & 28.1 $\to$ 1.5 & --- \\  % exp_id: 8b_cell_a_episode
Cell D: episode-level & \textbf{96.88} & \textbf{90.30} \\  % exp_id: 8b_cell_c_episode_v9d, 8b_v9d_ood_eval_v5
\bottomrule
\end{tabular}
\end{table}

Both the failure and the fix are scale-invariant. Step-independent training collapses within two steps at 8B, compared to 3--15 steps at 1.7B; episode-level training not only avoids collapse but exceeds the SFT ceiling by a substantial margin. Due to computational constraints of 8B-scale FSDP training, 8B results report single-seed runs.
